% =============================================================================
% Section 8: Discussion
% =============================================================================
\section{Discussion}\label{sec:discussion}

This section synthesises the empirical findings into five structural
conclusions and situates them within the broader context of strategic
autonomy measurement.

\subsection{Five structural conclusions}

\paragraph{Conclusion~1: The ISI is dominated by defence and logistics
variance.}
The variance decomposition (\cref{sec:variance-decomp}) shows that
the defence axis contributes 35.10\% and the logistics axis 34.26\%
of total composite variance, jointly accounting for 69.37\%.  The
own-variance component constitutes 64.28\% and the cross-covariance
component 35.72\% (\cref{tab:variance-decomp-detail}).  The practical
implication is that cross-country ranking differences are overwhelmingly
determined by two of the six axes, while the remaining four axes
(financial, energy, technology, critical inputs) together contribute
less than one-third of the composite's discriminatory power.

This dominance is not an artefact of scaling or weighting.  Because
all axes inhabit a common $[0,1]$ scale and receive equal weight,
the variance shares reflect genuine differences in cross-country
dispersion: defence has $\sigma = 0.226$ and logistics $\sigma = 0.248$,
while the financial axis has $\sigma = 0.037$.  The multiple
regression analysis (\cref{sec:axis-rank}) confirms this: defence and
logistics have the highest bivariate $R^2$ values (0.394 and 0.491,
respectively) and partial correlation coefficients (0.940 and 0.809)
with the composite, and a six-axis regression achieves $R^2 = 0.968$.

\paragraph{Conclusion~2: Rank ordering is robust except in the
compressed mid-tier.}
The robustness programme (\cref{sec:robustness}) demonstrates that
the ISI ranking is stable under winsorization ($\rho = 1.000$),
Dirichlet weight perturbation (mean $\rho = 0.979$), and five of
six LOO exercises ($\rho \geq 0.926$).  Only defence-axis removal
($\rho = 0.833$, $\tau = 0.675$) and z-score standardisation
($\rho = 0.896$, $\tau = 0.726$) produce moderate rank disruption.

The rank elasticity analysis (\cref{sec:rank-elasticity}) reveals that
13 of 27 countries are sensitive to $\pm 10\%$ axis perturbations, but
these are concentrated in the composite range 0.30--0.40 where
inter-country score gaps are smallest.  The extreme-ranked countries
(Malta, Cyprus, France, Germany) are entirely stable.  This pattern
is consistent with the compression diagnostic ($\eta^2 = 0.483$,
\cref{sec:compression-diagnostics}): the three-tier classification
explains only 48.3\% of within-composite variance, confirming that
the mid-range is genuinely compressed and inherently difficult to
resolve with \emph{any} aggregation method.

\paragraph{Conclusion~3: Classification compression is mathematically
inevitable.}
The compression diagnostics (\cref{tab:compression-diagnostics}) show
that the three-tier ISI classification (mildly, moderately, highly
concentrated) achieves $\eta^2 = 0.483$ using thresholds at 0.33 and
0.50.  Alternative classification schemes---tercile-based
($\eta^2 = 0.890$) and quartile-based ($\eta^2 = 0.933$)---achieve
higher between-group separation but sacrifice interpretability.

The low $\eta^2$ under the ISI thresholds reflects the composite
distribution itself: 17 of 27 countries fall in the ``moderately
concentrated'' tier (0.33--0.50), creating a crowded middle category
with within-tier dispersion of $\sigma = 0.052$.  This compression
is not a design flaw; it is a mathematical consequence of the
underlying data distribution, which is unimodal and moderately
right-skewed (\cref{fig:histogram}).  Any threshold-based
classification of a unimodal distribution will produce a crowded
middle tier.

\paragraph{Conclusion~4: European defence markets exhibit
oligopsonistic concentration.}
The defence structural analysis (\cref{sec:defence-structural})
reveals that 16 of 27 EU member states (59.3\%) have the defence
axis as their single highest-scoring axis.  Five countries exceed
0.80, and one country (Malta) reaches the theoretical maximum of
1.000 (\cref{tab:defence-structural}).  The mean gap between the
defence score and the next-highest axis is 0.184
(\cref{tab:defence-gap}).

This concentration is not an artefact of a few outliers: the
prevalence of defence dominance spans small (Malta, Cyprus),
medium (Netherlands, Bulgaria, Greece), and large (Italy) member
states.  The pattern is consistent with the oligopsonistic
structure of European defence procurement, in which a small number
of bilateral supplier relationships (predominantly US--EU) create
high HHI values that are difficult to diversify through conventional
trade policy.

\paragraph{Conclusion~5: Financial concentration is orthogonal to
goods-trade concentration.}
The correlation analysis (\cref{sec:cross-axis}) demonstrates that
the financial axis has near-zero Pearson correlations with all five
other axes ($|r| \leq 0.22$), with all $t$-statistics below 1.15
(\cref{tab:corr-tstat}).  The geographic correlates
(\cref{sec:geo-correlates}) confirm this: island status, population
size, and founding-member status---all strong predictors of
goods-trade concentration---have essentially no predictive power
for financial concentration.

This orthogonality has a structural interpretation.  The financial
axis measures the concentration of financial-service imports
(EBOPS classification), which depends on the structure of a country's
banking and insurance sectors rather than on its physical trade
geography.  A country can have highly concentrated goods imports
(high defence, logistics, energy scores) while maintaining
diversified financial-service imports, and vice versa.

The compensability analysis (\cref{sec:compensability}) further
illuminates this point: countries with high CV values (peaked
profiles) are overwhelmingly peaked on defence or logistics, not on
the financial axis.  The financial axis contributes only 1.03\% of
composite variance precisely because its cross-country dispersion
is so low ($\sigma = 0.037$).

\subsection{Methodological implications}

\paragraph{Equal weighting is defensible but not unique.}
The Dirichlet Monte Carlo exercise shows that equal weighting
produces rankings within the 95\% confidence envelope of randomly
weighted composites (mean $\rho = 0.979$).  This does not prove
that equal weighting is optimal---it proves that the ranking is
insensitive to moderate weight variation.  The eigenvalue analysis
(\cref{tab:eigenvalues}) supports this: no single principal
component explains more than 31.7\% of variance, so there is no
empirical basis for strongly favouring any particular axis.

\paragraph{Composite interpretation requires axis-level context.}
The compensability analysis demonstrates that the arithmetic-mean
composite can conceal peak vulnerabilities.  The Netherlands
(ISI = 0.335, defence = 0.960) and Estonia (ISI = 0.343, defence
= 0.470) have similar composites but radically different risk
profiles.  Policy recommendations derived solely from the composite
would be misleading; the axis-level profile must be consulted.

\paragraph{Threshold-based classification is inherently lossy.}
The compression diagnostics show that any threshold-based
classification of the ISI distribution will produce a crowded middle
tier.  Users of the ISI should treat the three-tier classification
as a coarse heuristic and consult the continuous composite and
axis-level scores for fine-grained comparisons.

\subsection{Limitations}

Several limitations constrain the interpretation of these results.

\begin{enumerate}[leftmargin=2em]
  \item \textbf{Cross-sectional design.}  The ISI~v1.0 is computed
    for a single year (reference year 2024).  Temporal dynamics---how
    concentration changes in response to policy interventions or
    geopolitical shocks---cannot be assessed without panel data.

  \item \textbf{HHI as the sole concentration measure.}  The axes
    use Herfindahl--Hirschman indices or comparable concentration
    statistics.  Alternative measures (entropy, Gini of trade shares,
    bilateral dependency ratios) might produce different rankings,
    particularly for countries near the median.

  \item \textbf{Equal weighting.}  Although the Dirichlet exercise
    shows rank stability under random weighting, there may be policy
    contexts in which defence concentration should receive higher
    weight than financial concentration.  The ISI framework
    accommodates differential weighting but does not prescribe it.

  \item \textbf{Data availability.}  Some axis scores rely on
    publicly available Eurostat and Comtrade data; others
    (particularly defence) may incorporate restricted or modelled
    data.  The reproducibility of individual axis scores depends on
    the transparency of the underlying data sources.

  \item \textbf{Small $N$.}  With $N = 27$, many pairwise correlations
    fail to reach statistical significance even when the point
    estimates are moderate (\cref{tab:corr-tstat}).  The geographic
    correlates (\cref{sec:geo-correlates}) are descriptive, not
    inferential, because the small sample precludes multivariate
    regression with more than a few covariates.
\end{enumerate}
